\section{Four Challenges Unique to Redistricting Inference}
\label{sec:challenges}

\subsection{Challenge C1: The Test Statistic Is a Plan}
\label{sec:challenge1}

In standard statistical testing, the test statistic is a scalar function
of the data (e.g., a sample mean, a regression coefficient).
Asymptotic distribution theory---central limit theorems, bootstrap
consistency results---applies to such scalars.

In redistricting, the object of inference is a plan: an assignment of
$n \approx 84{,}000$ census tracts to $k$ districts.
Summary statistics (compactness, seat counts, efficiency gap) are derived
from this assignment, but they do not capture the full geometry of the plan.
Two plans with identical compactness scores can be geometrically
completely different.

This creates a dimensionality problem: the ``parameter space'' of
redistricting plans is discrete and of dimension $O(n^k)$---astronomically
large.
Standard asymptotic theory does not apply.
Inference must proceed through simulation (the ensemble), and the quality
of inference depends entirely on the quality of the simulation.

\subsection{Challenge C2: The Null Distribution Requires Simulation}
\label{sec:challenge2}

The null distribution in redistricting testing is the distribution of
summary statistics (compactness, seat counts) across the space of all
valid plans---plans that are contiguous, population-balanced, and
otherwise legally compliant.

This distribution cannot be computed analytically.
The space of valid plans is too large to enumerate.
Instead, it is approximated by the ensemble: a set of plans drawn from
a Markov chain whose stationary distribution is uniform over valid plans.

The quality of this approximation depends on:
\begin{enumerate}
\item \textbf{Convergence}: Has the chain mixed? Does the empirical
  distribution of the ensemble approximate the true stationary distribution?
  G.4~\citep{g4paper} shows that convergence requires 5,000--50,000
  steps depending on the state.
\item \textbf{ESS}: What is the effective sample size, accounting for
  autocorrelation?
  A 10,000-step chain with lag-1 autocorrelation $\rho_1 = 0.87$
  (as in NC from G.4) has $\ess = \frac{10{,}000}{1 + 2 \cdot 0.87 /
  (1 - 0.87)} \approx 695$ independent samples---not 10,000.
\end{enumerate}

Using the nominal sample size (10,000) instead of the ESS (695) overstates
the precision of percentile estimates by a factor of $\sqrt{10{,}000 / 695}
\approx 3.8$.

\subsection{Challenge C3: Multiple Metrics Create a Testing Problem}
\label{sec:challenge3}

A standard redistricting analysis evaluates multiple metrics simultaneously:
Polsby-Popper compactness (or Reock), efficiency gap, mean-median, partisan
bias, the seats-votes curve slope, Democratic seat count, and potentially others.

If each metric is tested at level $\alpha = 0.05$, the family-wise error
rate (probability of at least one false positive) is:
\[
  \mathrm{FWER} = 1 - (1 - \alpha)^m \approx 1 - e^{-m\alpha},
\]
where $m$ is the number of tests.
For $m = 7$ metrics: $\mathrm{FWER} = 1 - (0.95)^7 \approx 0.30$.
Across 50 states and 3 cycles (1,050 tests): expected false positives
$= 1050 \times 0.05 = 52.5$ at the uncorrected $\alpha = 0.05$.

The standard practice of reporting ``the enacted plan is a significant
outlier on [metric X] at the 5\% level'' without adjustment for
multiple comparisons produces a substantial false-positive risk.

\subsection{Challenge C4: Ensembles Are Too Small for Tail Estimation}
\label{sec:challenge4}

The most extreme gerrymandering claims typically involve enacted plans
at the 1st or 99th percentile of the ensemble.
But tail percentiles have large sampling uncertainty.

Let $\hat{p}$ be the empirical percentile of the enacted plan in an
ensemble of size $n^*$ (ESS-corrected).
The standard error of $\hat{p}$ near the tails is approximately:
\[
  \mathrm{SE}(\hat{p}) \approx \sqrt{\frac{p(1-p)}{n^*}}.
\]
For $p = 0.01$ (1st percentile) and $n^* = 695$ (NC from G.4):
\[
  \mathrm{SE}(0.01) \approx \sqrt{\frac{0.01 \times 0.99}{695}}
  \approx 0.0038 = 0.38 \text{ percentage points}.
\]
The 95\% confidence interval on the true percentile is approximately
$[0.01 - 1.96 \times 0.38\%, \; 0.01 + 1.96 \times 0.38\%]
= [0.26\%, 1.74\%]$.

This interval is not wide for a 1\% claimed percentile.
But for a plan at the claimed 0.1st percentile with the same ESS,
the confidence interval includes 0\%---meaning the plan might not
actually be an outlier in the true stationary distribution.

S.3~\citep{s3paper} derives the minimum ensemble size for reliable
tail estimation at each significance level.
